% =========================================================================
%  Reporte — Actividad 3: Capa de visualización temática
%  Proyecto: Atlas Global de Desastres
%  Asignatura: Análisis y exploración de datos
% =========================================================================
\documentclass[11pt]{article}

\usepackage[utf8]{inputenc}
\usepackage[T1]{fontenc}
\usepackage[spanish,es-nodecimaldot]{babel}
\usepackage{geometry}
\usepackage{graphicx}
\usepackage{booktabs}
\usepackage{array}
\usepackage{xcolor}
\usepackage{hyperref}
\usepackage{listings}
\usepackage{caption}
\usepackage{float}

\geometry{margin=2.5cm}

\definecolor{azul}{HTML}{1D4ED8}
\definecolor{grisclaro}{HTML}{F1F5F9}
\definecolor{grisborde}{HTML}{CBD5E1}

\hypersetup{
  colorlinks=true,
  linkcolor=azul,
  urlcolor=azul,
  pdftitle={Actividad 3 - Capa de visualizacion tematica},
}

\lstset{
  basicstyle=\ttfamily\small,
  backgroundcolor=\color{grisclaro},
  frame=single,
  rulecolor=\color{grisborde},
  breaklines=true,
  showstringspaces=false,
}

\title{\textbf{Capa de visualización temática sobre el mapa 3D}\\[0.3em]
       \large Actividad 3 --- Atlas Global de Desastres}
\author{Instituto Tecnológico Superior de los Ríos\\
        Ingeniería en Sistemas Computacionales\\
        \small Asignatura: Análisis y exploración de datos\\
        \small Profesor: PhD. Fernando Pech-May}
\date{\today}

\begin{document}
\maketitle

\begin{abstract}
\noindent
Este reporte documenta el diseño e implementación de la capa de
visualización temática de la Actividad 3. La capa se construye sobre un
mapa mundial 3D y representa las variables de desastres naturales mediante
cuatro recursos visuales: coloración, intensidad de color, elevación y
deformación de la superficie. La aplicación permite seleccionar
simultáneamente dos variables ---una para el color y otra para la altura---
y alternar la proyección entre globo y plano. Se procesaron
10\,907 eventos EM-DAT sobre 175 unidades cartográficas.
\end{abstract}

\tableofcontents
\newpage

% -------------------------------------------------------------------------
\section{Introducción}
% -------------------------------------------------------------------------
La Actividad 3 del proyecto consiste en diseñar e implementar la capa de
visualización temática sobre el mapa 3D desarrollado en la Actividad 2,
incorporando recursos visuales como coloración, intensidad de color,
elevación y deformación de la superficie para representar las variables
seleccionadas.

\subsection*{Objetivo}
Representar de forma simultánea una o dos variables relacionadas con los
desastres naturales, de modo que a mayor valor de la variable mayor sea su
representación visual sobre el mapa.

% -------------------------------------------------------------------------
\section{Recursos visuales implementados}
% -------------------------------------------------------------------------
\begin{table}[H]
\centering
\caption{Correspondencia entre los recursos solicitados y su implementación.}
\begin{tabular}{@{}p{3.2cm}p{6.2cm}p{4.2cm}@{}}
\toprule
\textbf{Recurso} & \textbf{Implementación} & \textbf{Archivo} \\
\midrule
Coloración & Rampa de color interpolada en 12 paradas por variable &
\texttt{capa\_tematica.js} \\
Intensidad de color & Mezcla de cada color con su luminancia según un
control deslizante & \texttt{escalas.js} \\
Elevación & Capa \texttt{fill-extrusion}: prisma con altura proporcional al
valor & \texttt{capa\_tematica.js} \\
Deformación & \texttt{fill-extrusion} en relieve continuo de menor altura &
\texttt{capa\_tematica.js} \\
\bottomrule
\end{tabular}
\end{table}

\subsection{Coloración e intensidad}
Cada variable temática cuenta con una rampa perceptual propia
(\texttt{YlOrRd}, \texttt{Reds}, \texttt{Blues}, \texttt{Viridis},
\texttt{Magma}, \texttt{Teal}). Los territorios sin indicadores se pintan
con un color neutro para diferenciarlos de un valor cero. El control de
intensidad ajusta la saturación de la escala y permite regular la carga
visual del mapa.

\subsection{Elevación y deformación de la superficie}
La elevación se implementa con la capa \texttt{fill-extrusion} de MapLibre
GL, que proyecta cada polígono como un prisma sobre el globo. La altura en
metros se calcula como:
\[
h = h_{base} + t \cdot h_{max}, \qquad t \in [0, 1]
\]
donde $t$ es el valor normalizado de la variable de altura. La deformación
utiliza la misma capa con una altura máxima reducida y una base mayor,
generando un relieve continuo en lugar de prismas altos.

% -------------------------------------------------------------------------
\section{Arquitectura}
% -------------------------------------------------------------------------
El pipeline de datos se implementa en Python y produce una capa GeoJSON y
un archivo de metadatos que consume la aplicación web.

\begin{lstlisting}[language=Python, caption={Generación de la capa temática.}]
python -m src.capa_tematica
# -> app/datos/paises.geojson
# -> app/datos/metadatos.json
# -> resultados/tablas/resumen_capa_tematica.csv
\end{lstlisting}

\subsection{Tecnologías}
\begin{itemize}
  \item \textbf{Frontend:} HTML5, CSS3 y JavaScript vanilla (ES6).
  \item \textbf{Visualización 3D:} MapLibre GL JS (globo, elevación y relieve).
  \item \textbf{Backend:} Python, pandas, GeoPandas, NumPy.
  \item \textbf{Formato de capa:} GeoJSON.
\end{itemize}

% -------------------------------------------------------------------------
\section{Resultados}
% -------------------------------------------------------------------------
\begin{table}[H]
\centering
\caption{Cifras de la capa temática generada.}
\begin{tabular}{@{}lr@{}}
\toprule
\textbf{Métrica} & \textbf{Valor} \\
\midrule
Eventos EM-DAT procesados & 10\,907 \\
Unidades cartográficas & 175 \\
Territorios con indicadores & 168 \\
Periodo analizado & 2000--2026 \\
Tamaño del GeoJSON & $\approx$ 0.5 MB \\
\bottomrule
\end{tabular}
\end{table}

% Las capturas se generan con: python -m src.capturas
\begin{figure}[H]
  \centering
  \includegraphics[width=\textwidth]{imagenes/01_globo_eventos.png}
  \caption{Mapa de calor global por número de desastres (coloración).}
\end{figure}

\begin{figure}[H]
  \centering
  \includegraphics[width=\textwidth]{imagenes/02_globo_elevacion.png}
  \caption{Elevación de los territorios según los fallecidos.}
\end{figure}

\begin{figure}[H]
  \centering
  \includegraphics[width=\textwidth]{imagenes/03_globo_deformacion.png}
  \caption{Deformación de la superficie según personas afectadas.}
\end{figure}

\begin{figure}[H]
  \centering
  \includegraphics[width=\textwidth]{imagenes/05_combinado_eventos_fallecidos.png}
  \caption{Representación simultánea: color = eventos, elevación = fallecidos.}
\end{figure}

% -------------------------------------------------------------------------
\section{Conclusiones}
% -------------------------------------------------------------------------
La capa temática cumple con los recursos solicitados y permite comparar
simultáneamente dos variables. La separación entre el pipeline de datos
(Python) y la interfaz (JavaScript vanilla) facilita el mantenimiento y la
evolución del proyecto hacia las actividades 4 a 6.

\end{document}
